%%%%%%%%%%  scan idx 366  |  folio 353  %%%%%%%%%%
\(x-\gamma^n(x')\in(\Ch A)_{n+1}\). En
itérant l'opération \(p-n\) fois, on arrive à \((\Ch A)_{p+1}\), qui est nul, et on a donc écrit \(x\) comme un élément de \(\Fil^n\Ch A\). Compte tenu de (6.6.5), les deux filtrations sont donc identiques.

\textbf{6.7.} On suppose maintenant que \(A\) est un \(\lambda\)-anneau augmenté vers \(K\), l'homomorphisme d'augmentation étant noté \(\epsilon\). On dispose donc sur \(A\) de la filtration canonique de \(\lambda\)-anneau augmenté, et on peut donc associer à \(\Gr A\) son anneau de Chern.

Si \(x\in A\), alors \(\gamma_t(x-\epsilon(x))\in\Ghat(A)\) par définition de la filtration de \(\lambda\)-anneau augmenté. On note \(c(x)\) et on appelle \emph{classe de Chern totale de} \(x\) l'image de \(\gamma_t(x-\epsilon(x))\) dans \(\Ghat(A)/\Ghat'(A)=1+\widehat{\Gr A}^{+}\)~; on note \(c^i(x)\), et on appelle \(i\)-ième classe de Chern de \(x\), le \(i\)-ième coefficient de \(c(x)\)~; on a donc :
\begin{equation}\tag{6.7.1}
  c^i(x)=\cl_{\Gr^iA}\!\left(\gamma^i(x-\epsilon(x))\right)~;\qquad
  c(x)=1+c^1(x)+\cdots
\end{equation}
On pose alors :
\begin{equation}\tag{6.7.2}
  \widetilde c(x)=(\epsilon(x),c(x)),
\end{equation}
et on appelle \(\widetilde c(x)\) la classe de Chern complétée de \(x\).

\textbf{Proposition 6.8.} La classe de Chern complétée \(\widetilde c:A\longrightarrow\Ch(\Gr A)\), est un \(K\)-\(\lambda\)-homomorphisme de \(K\)-\(\lambda\)-algèbres augmentées.

Il est clair que l'homomorphisme ainsi défini est additif~; pour vérifier qu'il est multiplicatif, on écrit :
\[
\begin{aligned}
\gamma_t(xy-\epsilon(xy))
&=\gamma_t\bigl(\epsilon(y)(x-\epsilon(x))\bigr).
  \gamma_t\bigl(\epsilon(x)(y-\epsilon(y))\bigr).
  \gamma_t\bigl((x-\epsilon(x))(y-\epsilon(y))\bigr)\\
&=\bigl(\gamma_t(x-\epsilon(x))\bigr)^{\epsilon(y)}.
  \bigl(\gamma_t(y-\epsilon(y))\bigr)^{\epsilon(x)}.
  \bigl(\gamma_t(x-\epsilon(x))\nabla\gamma_t(y-\epsilon(y))\bigr)
\end{aligned}
\]
